\chapter{Appendix}\label{chap:appendix}

\section{Proofs Chapter~\ref{chap:ar}}

\subsection{Proof of Proposition~\ref{prop:hausdorff_distance_convergence_speed}}\label{sec:proof_hausdorff_distance_convergence_speed}

\begin{Lem}\label{lem:upper_bound_psi_j}
    There exist $C>0$ and $\rho\in(0,1)$ such that $|\psi_j|\le C\rho^{j}$. In particular, $\lim\limits_{j\to\infty}\psi_j=0$.
\end{Lem}

\begin{proof}
    Define $r_{\phi}:=\min\{|z|:z\in\mathbb{C},\phi(z)=0\}$. Since $\{X_t\}$ is causal, $r_{\phi}>1$. Let $R\in(1,r_{\phi})$ be arbitrary. From complex analysis, it follows that $1/\phi$ is analytical on the disk $D_0(R)=\{z\in\mathbb{C}:|z|\le R\}$. Let $1/\phi(z)=\sum_{j=0}^{\infty}\psi_jz^{j}$ be its representation. By Cauchy's integral inequality, see~\cite[Corollary~4.3]{Shiu01072004}, $|\psi_j|\le R^{-j}\max_{|z|=R}|1/\phi(z)|$. Setting $C=\max_{|z|=R}|1/\phi(z)|$ and $\rho=R^{-1}$ yields the statement.
\end{proof}

\begin{proof}
    First, notice that $\mathcal{Z}_n^{(k)}\subseteq\mathcal{Z}$ for all $n\in\mathbb{N}$ and therefore $d(y,\mathcal{Z}^{(k)})=0$ for all $y\in\mathcal{Z}_n^{(k)}$. By definition of the Hausdorff distance $d_H(\mathcal{Z}_n^{(k)},\mathcal{Z})=\sup\{d(\mathcal{Z}_n^{(k)},z):z\in\mathcal{Z}^{(k)}\}$. Let $y\in\mathcal{Z}_n^{(k)}$ and $z\in\mathcal{Z}^{(k)}$ be arbitrary. They have representations $\sum_{j=0}^{n}a_jg_j^{(k)}$ and $\sum_{j=0}^{\infty}b_jg_j^{(k)}$ respectively for some $a_j,b_j\in[-1,1]$. Now we have
    \begin{align*}
        d(y,z)&=\left\|\sum_{j=0}^{n}a_jg_j^{(k)}-\sum_{j=0}^{\infty}b_jg_j^{(k)}\right\|=\left\|\sum_{j=n+1}^{\infty}b_jg_j^{(k)}+\sum_{j=0}^{n}(b_j-a_j)g_j^{(k)}\right\|\\
        &\le\sum_{j=n+1}^{\infty}|b_j|\|g_j^{(k)}\|+\sum_{j=0}^{n}|b_j-a_j|\|g_j^{(k)}\|
    \end{align*}
    and consequently
    \begin{align*}
        d(\mathcal{Z}_n^{(k)},z)&=\inf_{a_0,\dots,a_n}d\left(\sum_{j=0}^{n}a_jg_j^{(k)},z\right)\\&\le\inf_{a_0,\dots,a_n}\left(\sum_{j=n+1}^{\infty}|b_j|\|g_j^{(k)}\|+\sum_{j=0}^{n}|b_j-a_j|\|g_j^{(k)}\|\right)\\&
        =\sum_{j=n+1}^{\infty}|b_j|\|g_j^{(k)}\|
    \end{align*}
    Finally,
    \begin{align*}
        d_H(\mathcal{Z}_n^{(k)},\mathcal{Z})=\sup_{b_0,b_1,\dots}d\left(\mathcal{\mathcal{Z}_n^{(k)}},\sum_{j=0}^{\infty}b_jg_j^{(k)}\right)\le\sup_{b_0,b_1,\dots}\sum_{j=n+1}^{\infty}|b_j|\|g_j^{(k)}\|=\sum_{j=n+1}^{\infty}\|g_j^{(k)}\|
    \end{align*}
    Using Lemma~\ref{lem:upper_bound_psi_j} we obtain
    \begin{align*}
        \|g_j^{(k)}\|=\sqrt{\sum_{r=0}^{k-1}|\psi_{j-r}|^2}\le\sqrt{\sum_{r=0}^{k-1}C^2\rho^{2j}\rho^{-2r}}=C\rho^{j-k+1}\sqrt{\frac{1-\rho^{2k}}{1-\rho^2}}\le C\sqrt{k}\rho^{j-k+1}
    \end{align*}
    In the end
    \begin{align*}
        d_H(\mathcal{Z}_n^{(k)},\mathcal{Z}^{(k)})\le\sum_{j=n+1}^{\infty}\|g_j^{(k)}\|\le C\sqrt{k}\frac{1}{\rho^{k-2}}\rho^{n}\sum_{j=0}^{\infty}\rho^{j}\le C\frac{\sqrt{k}}{1-\rho}\frac{1}{\rho^{k-2}}\rho^{n}
    \end{align*}
\end{proof}

\subsection{Proof of Proposition~\ref{prop:minimum_width}}\label{sec:proof_minimum_width}

\begin{Lem}\label{lem:young_convolution_equality}
    Let $a=(a_n)_{n\in\mathbb{Z}}\in l^2(\mathbb{Z})\setminus\{0\}$ and $b=(b_n)_{n\in\mathbb{Z}}\in l^1(\mathbb{Z})$ be sequences. Then it holds $\|a\ast b\|_2=\|a\|_2\|b\|_1$ if and only if $b=0$ or $b=c\delta_r$ for some $c\ne0$ and $r\in\mathbb{Z}$.
\end{Lem}

\begin{proof}
    First, assume that the equality holds. We define the shifted sequence $\tau_ma$ by $(\tau_ma)[n]=a_{n-m}$. Then the convolution can be seen as an infinitely weighted sum of shifts, that is, $a\ast b=\sum_mb_m\tau_ma$. We set $x_m=b_m\tau_ma$. By assumption $\{x_m\}$ is absolutely summable, which implies $\|\sum_{|m|\le N}x_m-\sum_mx_m\|\to0$ for $N\to\infty$ and by the reverse triangle inequality $|\|\sum_{|m|\le N}x_m\|_2-\|\sum_mx_m\|_2|\le\|\sum_{|m|\le N}x_m-\sum_mx_m\|_2$. Therefore,
    \begin{align*}
        \sum_{m,k}\langle x_m,x_k\rangle=\lim_{N\to\infty}\sum_{|m|\le N}\sum_{|k|\le N}\langle x_m,x_k\rangle=\lim_{N\to\infty}\left\|\sum_{|m|\le N}x_m\right\|_2^2=\left\|\sum_mx_m\right\|_2^2.
    \end{align*}
    By assumption we also have
    \begin{align*}
        \sum_m\|x_m\|_2=\|a\|_2\|b\|_1=\|a\ast b\|_2=\left\|\sum_mx_m\right\|_2
    \end{align*}
    Thus,
    \begin{align*}
        \sum_{m,k}\langle x_m,x_k\rangle=\left\|\sum_{m}x_m\right\|_2^2=\sum_{m,k}\|x_m\|_2\|x_k\|_2
    \end{align*}
    By Cauchy Schwarz, $\langle x_m,x_k\rangle\le\|x_m\|_2\|x_k\|_2$ and because of the above equality it follows $\langle x_m,x_k\rangle=\|x_m\|_2\|x_k\|_2$. This means that $x_m$ and $x_k$ are linearly dependent for all $m,k$ with $b_m,b_k\ne0$ and therefore $x_m=c_{m,k}x_k$ for some $c_{m,k}\ne0$. It further follows that $|a_{n-m}|=|a_{n-k}|$ for all $n$. Now $m\ne k$ would lead to $a\notin l^{2}(\mathbb{Z})$. This shows that $b$ must either be a null sequence or can only have one nonzero entry.

    The other direction simply follows by plugging in and checking.
\end{proof}

\begin{proof}
    Recall that $g_j^{(k)}=(\psi_j,\dots,\psi_{j-k+1})^{T}$ where $\{\psi_{j}\}_{j\in\mathbb{Z}}$ is a sequence given by
    \begin{align*}
        \psi_j=\begin{cases}
            \sum_{i=1}^{p}\phi_i\psi_{j-i} & ,j\ge 1\\
            1 & ,j=0\\
            0 & ,j<0.
        \end{cases}
    \end{align*}
    Let $u=(u_0,\dots,u_{k-1})\in\mathbb{S}^{k-1}$ and consider the extension of $u$ to a sequence $\{\tilde{u}_l\}_{l\in\mathbb{Z}}$ where $\tilde{u}_l=\mathbf{1}\{0\le l\le k-1\}u_l$. Define the sequence $a=\{a_j\}_{j\in\mathbb{Z}}$ by
    \begin{align*}
        a_j=\langle g_j^{(k)},u\rangle=\sum_{l=0}^{k-1}u_l\psi_{j-l}=\sum_{l\in\mathbb{Z}}\tilde{u}_l\psi_{j-l}=(\tilde{u}\ast\psi)[j].
    \end{align*}
    By Definition of $\{\psi_j\}_{j\in\mathbb{Z}}$ we have $(q_{\phi}\ast\psi)=\delta_0$ and thus $q_{\phi}\ast a=q_{\phi}\ast(\tilde{u}\ast\psi)=(q_{\phi}\ast\psi)\ast\tilde{u}=\delta_0\ast\tilde{u}=\tilde{u}$ by properties of the convolution. By Young's convolution inequality (see~\cite[Theorem 1.2.12]{grafakos_classical_2014})
    \begin{align}\label{eq:young_convolution_inequality}
        1=\|u\|_2=\|q_{\phi}\ast a\|_2\le\|q_{\phi}\|_2\|a\|_1.
    \end{align}
    Next, we use the simple formula for the support function of the zonoid $\mathcal{Z}^{(k)}_0$ to obtain
    \begin{align*}
        w_k(u)=(b-a)h_{\mathcal{Z}^{(k)}_0}(u)=(b-a)\sum_{j=0}^{\infty}|\langle g_j^{(k)},u\rangle|=(b-a)\|a\|_1.
    \end{align*}
    Therefore, we obtain
    \begin{align*}
        w_k(u)=(b-a)\|a\|_1\ge\frac{b-a}{\|q_{\phi}\|_2}.
    \end{align*}
    Plugging in the directions from the set in Equation~\eqref{eq:minimizers} in $w_k$ we easily see that the above lower bound is actually attained. It remains to show that there are no other minimizing directions. Inequality~\eqref{eq:young_convolution_inequality} becomes an equality iff $a$ has exactly one nonzero entry (see Lemma~\ref{lem:young_convolution_equality} in the Appendix~\ref{chap:appendix}). This means $a=c\delta_r$ and it follows $\tilde{u}=c\tau_rq_{\phi}$, which forces $u$ to have the form $\pm q_r^{(k)}/\|q_{\phi}\|_2$ for some $r=0,\dots,k-(p+1)$.

    Now let $k<p+1$. We have shown above that $w_k(u)$ attains the lower bound $(b-a)/\|q_{\phi}\|_2$ if and only if $u$ has the specific form above. But if $k<p+1$ there are no permitted $r$ that would lead to the specific form. Since $w_k$ is continuous it attains its minimum on $\mathbb{S}^{k-1}$ and therefore $m_k>(b-a)/\|q_{\phi}\|_2$.

    We are left to show $m_{k-1}\ge m_k$ for $k=2,\dots,p$. Let $v=(v_0,\dots,v_{k-2})\in\mathbb{S}^{k-2}$ and $u=(v_0,\dots,v_{k-2},0)\in\mathbb{S}^{k-1}$ and thus
    \begin{align*}
        w_{k}(u)=(b-a)\sum_{j}\left|\sum_{l=0}^{k-1}u_l\psi_{j-l}\right|=(b-a)\sum_{j}\left|\sum_{l=0}^{k-2}v_l\psi_{j-l}\right|=w_{k-1}(v).
    \end{align*}
    Let $v$ be a minimizing direction. Then $m_{k-1}=w_{k-1}(v)=w_k(u)\ge\min_{u\in\mathbb{S}^{k-1}}w_k(u)=m_k$.
\end{proof}

\subsection{Proof of Lemma~\ref{lem:hausdorff_convergence}}\label{sec:proof_hausdorff_convergence}

\begin{Lem}\label{lem:property_hausdorff_distance}
    Let $A,B,C$ be subsets of a metric space with $A\subseteq B\subseteq C$. Then $d_H(B,C)\le d_H(A,C)$.
\end{Lem}

\begin{proof}
    Since $A\subseteq B$, we have $\inf_{b\in B}d(b,c)\le\inf_{a\in A}d(a,c)$. Therefore,
    \begin{align*}
        d_H(B,C)=\sup_{c\in C}\inf_{b\in B}d(b,c)\le\sup_{c\in C}\inf_{a\in A}d(a,c)=d_H(A,C).
    \end{align*}
\end{proof}

\begin{proof}
    Let $\epsilon>0$ be arbitrary. Let $\Omega_0=\{\omega\in\Omega:\overline{\{Y_t(\omega):t\ge1\}}=\mathcal{Z}\}$. By assumption, $\mathbb{P}(\Omega_0)=1$. Fix $\omega\in\Omega_0$. In what follows, we abbreviate $Y_t=Y_t(\omega)$ and $\mathcal{P}_T=\{Y_1(\omega),\dots,Y_T(\omega)\}$. Because $\{Y_t:t\ge1\}$ is dense in $\mathcal{Z}$, it holds $\mathcal{Z}\subseteq\bigcup_{t\ge1}B(Y_t,\epsilon)$ and since $\mathcal{Z}$ is compact, there exists a finite sub-cover $\mathcal{Z}\subseteq\bigcup_{j=1}^{m(\epsilon)}B(Y_{t_j(\epsilon)},\epsilon)$. Set $T_{\epsilon}=\max\{t_1(\epsilon),\dots,t_{m(\epsilon)}(\epsilon)\}$. Now, let $z\in\mathcal{Z}$ be arbitrary. It follows that $d(z,\mathcal{P}_T)=\inf\{\|z-Y_t\|:t=1,\dots,T\}\le\epsilon$ for all $T\ge T_{\epsilon}$. Thus, $\sup_{z\in\mathcal{Z}}d(z,\mathcal{P}_T)\le\epsilon$, since $z$ was arbitrary. Since $\mathcal{P}_T\subseteq\mathcal{Z}$, we also have $\sup_{t=1,\dots,T}d(Y_t,\mathcal{Z})=0$. By definition of the Hausdorff distance, we obtain $d_H(\mathcal{P}_T,\mathcal{Z})=\sup_{z\in\mathcal{Z}}d(z,\mathcal{P}_T)\le\epsilon$ and therefore $d_H(\mathcal{P}_T,\mathcal{Z})\to0$ for $T\to\infty$, since $\epsilon>0$ was arbitrary. By Lemma~\ref{lem:property_hausdorff_distance} we also get $d_H(\operatorname{conv}(\mathcal{P}_T),\mathcal{Z})\to0$ for $T\to\infty$. Since we worked on $\Omega_0$, this proves almost sure convergence.
\end{proof}

\subsection{Proof of Lemma~\ref{lem:ergodic_implies_dense}}\label{sec:proof_ergodic_implies_dense}

\begin{proof}
    It is well known that the family $\mathcal{B}$ of open balls with rational centers $q\in\mathbb{Q}^{k}$ and radii $r\in\mathbb{Q}_{>0}$ forms a countable base for $\mathbb{R}^{k}$. That is, $\mathcal{B}$ is countable and for any open set $U\subseteq\mathbb{R}^k$ and every $x\in U$ there exists $B\in\mathcal{B}$ such that $x\in B\subseteq U$. Define $\mathcal{B}_{\mathcal{Z}}=\{B\in\mathcal{B}:B\cap\mathcal{Z}\ne\emptyset\}$. Let $B\in\mathcal{B}_{\mathcal{Z}}$ and choose some $z\in B$. Because $B$ is open, there exists $\delta>0$ with $B(z,\delta)\subseteq B$. Since $z\in\mathcal{Z}=\operatorname{supp}(Y_t)$, by the definition of the support, we have $\mathbb{P}(Y_t\in B)\ge\mathbb{P}(Y_t\in B(z,\delta))>0$. By Birkhoff's ergodic theorem
    \begin{align*}
        \frac{1}{T}\sum_{t=1}^{T}1\{Y_t\in B\}\stackrel{T\to\infty}{\to}\mathbb{P}(Y_t\in B)>0
    \end{align*}
    This implies that $Y_t\in B$ for infinitely many $t\ge 1$ almost surely. In particular, there exists $t\ge1$ such that $Y_t\in B$. Let $z\in\mathcal{Z}$ and $\epsilon>0$ be arbitrary. There exists $B\in\mathcal{B}$ such that $z\in B\subseteq B(z,\epsilon)$. From above we also know that there exists $t\ge1$ such that $Y_t\in B$. Thus $\|Y_t-z\|\le\epsilon$
\end{proof}

\subsection{Proof of Proposition~\ref{prop:sliding_window_fails_brunel}}\label{sec:proof_sliding_window_fails_brunel}

\todo[inline]{Extend to Beta distributed innovations}

\begin{proof}
    Let $\{Z_t\}\stackrel{\text{iid}}{\sim}U[-1,1]$ and write $X_t=\sum_{j=0}^{\infty}\psi_jZ_{t-j}$. Then, $\operatorname{supp}(X_t)=[-H,H]$, where $H=\sum_{j=0}^{\infty}|\psi_j|$. Consequently
    \begin{align*}
        H-X_t=\sum_{j=0}^{\infty}|\psi_j|-\sum_{j=0}^{\infty}\psi_jZ_{t-j}=\sum_{j=0}^{\infty}|\psi_j|\underbrace{(1-\operatorname{sgn}(\psi_j)Z_{t-j})}_{=D_{t-j}\in[0,2]}.
    \end{align*}
    There exist infinitely many nonzero $\psi_j$. (Assume, for contradiction, that only finitely many $\psi_j$ are nonzero. Then, the set $\{j\in\mathbb{N}_0:\psi_j\ne0\}$ is nonempty ($\psi_0=1$) and finite. Let $J$ be its maximum. This would mean that $1/\phi(z)=\sum_{j=0}^{J}\psi_jz^{j}$ is a polynomial of order $J$. By assumption, the polynomial $\phi$ has degree $p$. But now we have $0=\deg(\phi\cdot(1/\phi))=\deg(\phi)+\deg(1/\phi)=p+J$. This would force $p=0$ and $J=0$, which contradicts our assumption that $p\ge1$.) Thus, for each $\alpha>0$ no matter how large we find distinct indices $j_1,\dots,j_M$ where $M>\alpha$ such that $\psi_{j_m}\ne0$ for all $m=1,\dots,M$. Let $\epsilon>0$. Then $H-X_t\le\epsilon$ implies $D_{t-j_m}\le\epsilon|\psi_{j_m}|^{-1}$ for all $m=1,\dots,M$. Now, it follows
    \begin{align*}
        \mathbb{P}(H-X_t\le\epsilon)&\le\mathbb{P}\left(\bigcap_{m=1}^{M}\left\{D_{t-j_m}\le\epsilon|\psi_{j_m}|^{-1}\right\}\right)\stackrel{(\ast)}{=}\prod_{m=1}^{M}\mathbb{P}(D_{t-j_m}\le\epsilon|\psi_{j_m}|^{-1})\\&
        \stackrel{(\ast\ast)}{=}\prod_{m=1}^{M}\min\left\{\frac{\epsilon}{2|\psi_{j_m}|},1\right\}\le\prod_{m=1}^{M}\frac{\epsilon}{2|\psi_{j_m}|}=C_M\epsilon^M,
    \end{align*}
    where in step $(\ast)$ we have used the independence of the innovations and in step $(\ast\ast)$ that $D_{t-j_m}\sim\mathcal{U}[0,2]$. Because $H-X_t\le\epsilon\iff Y_t^{(k)}\in C_{Z^{(k)}}(e_1,\epsilon)$, we conclude
    \begin{align*}
        \frac{\mathbb{P}(Y_t^{(k)}\in C_{\mathcal{Z}^{(k)}}(e_1,\epsilon))}{\epsilon^{\alpha}}=\frac{\mathbb{P}(H-X_t\le\epsilon)}{\epsilon^{\alpha}}\le C_M\epsilon^{M-\alpha}\stackrel{\epsilon\to0}{\to}0.
    \end{align*}
\end{proof}

\subsection{Proof of Proposition~\ref{prop:properties_empirical_width}}\label{sec:proof_properties_empirical_width}

\begin{proof}
    Let $u,v\in S^{k-1}$. Write $\hat{w}_k(u)=\max_{s,t}\langle Y_s^{(k)}-Y_t^{(k)},u\rangle$ and we obtain
    \begin{align*}
        \hat{w}_k(u)-\hat{w}_k(v)&=\max_{s,t}\langle Y_s^{(k)}-Y_t^{(k)},u\rangle-\max_{s,t}\langle Y_s^{(k)}-Y_t^{(k)},v\rangle\\
        &\le\max_{s,t}\left(\langle Y_s^{(k)}-Y_t^{(k)},u\rangle-\langle Y_s^{(k)}-Y_t^{(k)},v\rangle\right)\\
        &=\max_{s,t}\langle Y_s^{(k)}-Y_t^{(k)},u-v\rangle\\
        &\le\left(\max_{s,t}\|Y_s^{(k)}-Y_t^{(k)}\|\right)\|u-v\|
    \end{align*}
    for all $u,v\in\mathbb{R}^{k}$.

    Convexity of $\hat{w}_k$ follows from the convexity of linear functions and the pointwise maximum of convex functions. Piecewise linearity also follows directly from the standard fact that the pointwise maximum of linear functions is piecewise linear. Differentiability almost everywhere follows directly from Rademacher's theorem (Lipschitz-continuous functions are differentiable almost everywhere). The set of non-differentiable points is just the set of all points $u$ where two distinct linear functions $\langle Y_s^{(k)}-Y_t^{(k)},u\rangle$ and $\langle Y_q^{(k)}-Y_p^{(k)},u\rangle$ coincide. The gradient of $\hat{w}_k$ in all directions $u\notin\mathcal{N}$ is obtained by straightforward differentiation.

    By Weierstrass' extreme value theorem, $\hat{w}_k$ has a global minimum on $S^{k-1}$. The set of minimizers
    \begin{align}\label{eq:set_of_minimizers}
        \mathcal{X}_k=\arg\min\{\hat{w}_k(u):u\in\mathbb{S}^{k-1}\}
    \end{align}
    is not empty. We also note that $\mathcal{X}_k$ is compact, again due to the continuity of $\hat{w}_k$ and the compactness of $\mathbb{S}^{k-1}$. 
    
    Denote by $\hat{m}_{k,T}$ the minimum of $\hat{w}_k$ on $\mathbb{S}^{k-1}$. We have $\hat{m}_{k,T}>0$ if and only if the affine hull of $\mathcal{P}_k$ is the whole space $\mathbb{R}^{k}$. Let $\operatorname{aff}(\mathcal{P}_{k,T})=\mathbb{R}^{k}$. Assume for a contradiction that $\hat{m}_{k,T}=0$. Then there would exist a direction $u\in\mathbb{S}^{k-1}$ with $\langle Y_t^{(k)},u\rangle=c_{k,u}$ constant for all observations $t$. Hence, all points $Y_t^{(k)}$ lie inside the affine hyperplane $H_u=\{x\in\mathbb{R}^{k}:\langle u,x\rangle=c_{k,u}\}$ and therefore $\operatorname{aff}(\mathcal{P}_{k,T})\subseteq H_u$. But $H_u$ is not the whole space, leading to a contradiction. If $\operatorname{aff}(\mathcal{P}_{k,T})\ne\mathbb{R}^{k}$, then $V:=\{Y_t^{(k)}-Y_K^{(k)}:t=K,\dots,T\}$ is a proper subspace since $Y_K^{(k)}+V=\operatorname{aff}(\mathcal{P}_{k,T})$. This means that we can find a direction $u'\in V^{\perp},\|u'\|=1$ and this direction fulfills $\langle Y_t^{(k)}-Y_K^{(k)},u\rangle=0$ for all $t=K,\dots,T$, implying that $\langle Y_t^{(k)},u\rangle$ is constant for all $t$ and thus $\hat{m}_{k,T}=0$.

    Take $Y_K^{(k)},\dots,Y_{K+k}^{(k)}$ and let the matrix $M$ be defined by having columns $Y_{K+1}^{(k)}-Y_K^{(k)},\dots,Y_{K+k}^{(k)}-Y_{K}^{(k)}$. The determinant $D=\det{M}$ is a nonzero polynomial in the observations $W=(X_{K-k+1},\dots,X_K,X_{K+1},\dots,X_{K+k})$, implying that $\{x:D(x)=0\}$ has Lebesgue measure zero; see~\cite{caron_zero_nodate}. Note that we need at least $T\ge K+k$ to form this vector $W$. The random vector $W$ has a joint density if the innovations have a density. Therefore, it follows $\mathbb{P}(D(W)\ne0)=1$, that is, the determinant of $M$ is almost surely nonzero. This immediately implies that the differences $Y_{K+1}^{(k)}-Y_K^{(k)},\dots,Y_{K+k}^{(k)}-Y_{K}^{(k)}$ span the whole space $\mathbb{R}^{k}$ and thus $\operatorname{aff}(\mathcal{P}_{k,T})=\mathbb{R}^{k}$. Combining the arguments, we find that $\hat{m}_{k,T}>0$ is almost surely the case given that $T\ge K+k$.
\end{proof}

\subsection{Proof of Lemma~\ref{lem:weak_sharp_minima_application}}\label{sec:proof_weak_sharp_minima_application}

\begin{proof}
    The function $\hat{w}_k$ is a pointwise maximum over finitely many linear functions. In the following, we write $A=\{Y_s^{(k)}-Y_t^{(k)}:s,t\}$ and $f(u)=\max\{\langle a,u\rangle:a\in A\}$. Also, let $m=\min\{f(u):u\in\mathbb{S}^{k-1}\}$ and $\mathcal{X}=\{u\in\mathbb{S}^{k-1}:f(u)=m\}$. Let $u^{\ast}\in\mathcal{X}$ and $T_{u^{\ast}}\mathbb{S}^{k-1}$ be the tangent space at $u^{\ast}$. Furthermore, let $A_{\ast}=\{a\in A:\langle a,u^{\ast}\rangle=m\}$ be the set of active vectors. Define $g(d)=\max\{\langle a,d\rangle:a\in A_{\ast}\}$ for all unit tangent directions $d\in T_{u^{\ast}}\mathbb{S}^{k-1}\cap\mathbb{S}^{k-1}=:S_{u^{\ast}}\mathbb{S}^{k-1}$.

    Let $t>0$. We can see that
    \begin{align*}
        f\left(u^{\ast}+td\right)=\max_{a\in A}\langle a,u^{\ast}+td\rangle=\underbrace{\max_{a\in A}\langle a,u^{\ast}\rangle}_{=m}+t\max_{a\in A}\langle a,d\rangle.
    \end{align*}
    We want to assert that $f(u^{\ast}+td)=m+tg(d)$ for sufficiently small $t>0$. In general, $\max\{\langle a,d\rangle:a\in A\}\ge\max\{\langle a,d\rangle:a\in A_{\ast}\}=g(d)$ since $A_{\ast}\subseteq A$. In case of equality there is nothing to show. Therefore we assume $\max\{\langle a,d\rangle:a\in A\}>g(d)$. If this is the case, there must exist an inactive vector $b\in A\setminus A_{\ast}$ with $\langle b,d\rangle>g(d)$. We now write
    \begin{align*}
        \langle b,u^{\ast}+td\rangle=\langle b,u^{\ast}\rangle+t\langle b,d\rangle=m-(\underbrace{m-\langle b,u^{\ast}\rangle}_{=:\delta_b>0})+t\langle b,d\rangle.
    \end{align*}
    We see that the difference
    \begin{align*}
        m+tg(d)-\langle b,u^{\ast}+td\rangle=\delta_b-t\left(\langle b,d\rangle-g(d)\right)
    \end{align*}
    stays positive for $t<\delta_b/(\langle b,d\rangle-g(d))=:r_b>0$. Set $r(u^{\ast})$ to be the minimum over all $b\in A\setminus A_{\ast}$ with $\langle b,d\rangle>g(d)$. It follows that
    \begin{align*}
        f(u^{\ast}+td)=\max_{a\in A}\langle a,u^{\ast}+td\rangle=\max_{a\in A_{\ast}}\langle a,u^{\ast}+td\rangle=m+tg(d)
    \end{align*}
    for all $t<r(u^{\ast})$ and $d\in S_{u^{\ast}}\mathbb{S}^{k-1}$ and therefore also
    \begin{align}\label{eq:proof_weak_sharp_minima_lemma01}
        f\left(\frac{u^{\ast}+td}{\sqrt{1+t^2}}\right)=\frac{m+tg(d)}{\sqrt{1+t^2}}
    \end{align}
    for all $d\in S_{u^{\ast}}\mathbb{S}^{k-1}$ and all $0<t<r(u^{\ast})$.

    From this it follows that $g(d)>0$ and therefore also $c^{\ast}=\min\{g(d):d\in S_{u^{\ast}}\mathbb{S}^{k-1}\}>0$ since $g$ is a continuous function and the unit tangent sphere $S_{u^{\ast}}\mathbb{S}^{k-1}$ is compact. Now write $u_t(d)=(u^{\ast}+td)/\sqrt{1+t^2}$. For $t\in(0,r(u^{\ast}))$ we have
    \begin{align*}
        f(u_t(d))-m&\stackrel{\eqref{eq:proof_weak_sharp_minima_lemma01}}{=}\frac{m+tg(d)}{\sqrt{1+t^2}}-m\ge\frac{m+tc^{\ast}}{\sqrt{1+t^2}}-m\\&
        =tc^{\ast}+(m+tc^{\ast})\left(\frac{1}{\sqrt{1+t^2}}-1\right)\\&
        =tc^{\ast}-t^2\underbrace{\frac{m+tc^{\ast}}{\sqrt{1+t^2}(1+\sqrt{1+t^2})}}_{\to m/2\text{ for }t\downarrow0}.
    \end{align*}
    Denote the remainder term by $r_t$. We have $|r_t|t^{-2}\to m/2$ for $t\downarrow0$ and in particular there exist constants $C>0$ and $r_{\ast}'>0$ such that $|r_t|\le Ct^2$ for all $t\in(0,r_{\ast}')$.

    We continue working on $t<\min\{r(u^{\ast}),r_{\ast}',\sqrt{2},c^{\ast}/(2C)\}$. Because $\|u_t(d)-u^{\ast}\|^2=2-2/\sqrt{1+t^2}\le\min\{t^2,2\}=t^2$ it follows
    \begin{align}\label{eq:proof_weak_sharp_minima_lemma02}
        f(u_t(d))-m\ge tc^{\ast}+r(t)\ge tc^{\ast}-Ct^2=t(c^{\ast}-Ct)\ge\frac{c^{\ast}}{2}t\ge\frac{c^{\ast}}{2}\|u_t(d)-u^{\ast}\|,
    \end{align}
    which lets us see that $\mathcal{X}$ must be finite.

    To see this, assume that $u_t(d)\in\mathcal{X}$ is another minimizer. It would follow from Equation~\eqref{eq:proof_weak_sharp_minima_lemma01} that $f(u_t(d))-m\ge(c^{\ast}/2)\|u_t(d)-u^{\ast}\|>0$, which is a contradiction to $u_t(d)\in\mathcal{X}$. Thus for each $u^{\ast}\in\mathcal{X}$ there exists a radius $s(u^{\ast})=\min\{\dots\}$ such that $B(u^{\ast},s(u^{\ast}))\cap\mathcal{X}=\{u^{\ast}\}$. Obviously $\{B(u^{\ast},s(u^{\ast})):u^{\ast}\in\mathcal{X}\}$ covers $\mathcal{X}$. Since $\mathcal{X}$ is compact, there exists a finite subcover. It follows that $\mathcal{X}$ must be finite and we denote by $u^{\ast}_1,\dots,u^{\ast}_q\in\mathcal{X}$ the finitely many minimizing directions. We can thus choose $\alpha=\min_{1\le i\le q}c^{\ast}_i/2>0$ and $\rho=\min_{1\le i\le q}s(u^{\ast}_i)$. For all $u\in\mathbb{S}^{k-1}$, $\operatorname{dist}(u,\mathcal{X})<\rho$ now follows
    \begin{align*}
        f(u)-m\ge\|u-u_i^{\ast}\|\frac{c_i^{\ast}}{2}\ge\alpha\operatorname{dist}(u,\mathcal{X})
    \end{align*}
\end{proof}

\section{Proofs Chapter~\ref{chap:kde}}

\subsection{Proof of Lemma~\ref{lem:acute_chain_perron_vec}}\label{sec:proof_acute_chain_perron_vec}

\begin{proof}
    Since the points $x_j$ are distinct, all entries of the matrix $D$ are positive and by Perron-Frobenius, its Perron vector $v$ has positive entries and is unique up to scaling. We fix $v$ so that $v^{T}\mathbf{1}=1$. We interpret $v$ as a discrete probability distribution on $\{x_1,\dots,x_N\}$. Let $X$ be a discrete random variable with $\mathbb{P}(X=x_j)=v_j$. Write the expected value $\mathbb{E}[X]=\mu=\sum_{j=1}^{N}v_jx_j$ and the variance $\operatorname{Var}[X]=\sigma^2=\sum_{j=1}^{N}v_j\|x_j-\mu\|^2$. Now we have
    \begin{align}\label{eq:bias_variance_decomp_perron_entry}
        (Dv)_i=\sum_{j=1}^{N}D_{ij}v_j=\sum_{j=1}^{N}\|x_i-x_j\|^2v_j=\mathbb{E}[\|X-x_i\|^2]\stackrel{(\ast)}{=}\|\mu-x_i\|^2+\sigma^2,
    \end{align}
    where at $(\ast)$ we have used the bias-variance decomposition. Since $v$ is the Perron vector of $D$, the left hand side of Equation~\eqref{eq:bias_variance_decomp_perron_entry} is equal to $\rho v_i$ with $\rho>0$ being the unique maximum eigenvalue of $D$.

    Next, we write
    \begin{align*}
        \mu&=\sum_{j=1}^{N}v_jx_j=v_1x_1+v_2x_2+\dots+v_Nx_N\\&
        =x_1-(v_2+\dots+v_N)x_1\\&
        +v_2x_2+v_3x_2+\dots+v_Nx_2-(v_3x_2+\dots+v_Nx_2)\\&
        \vdots\\&
        +v_{N-2}x_{N-2}+v_{N-1}x_{N-2}+v_Nx_{N-2}-(v_{N-1}x_{N-2}+v_Nx_{N-2})\\&
        +v_{N-1}x_{N-1}+v_Nx_{N-1}-v_Nx_{N-1}\\&
        +v_Nx_N\\&
        =x_1+(v_2+\dots+v_N)(x_2-x_1)+\dots+(v_{N-1}+v_N)(x_{N-1}-x_{N-2})+v_N(x_N-x_{N-1})\\&
        =x_1+\sum_{k=1}^{N-1}\sum_{j=k+1}^{N}v_j(x_{k+1}-x_k)=x_1+\sum_{k=1}^{N-1}a_k\delta_k,
    \end{align*}
    where $a_k=\sum_{j=k+1}^{N}v_j$ and $\delta_k=x_{k+1}-x_k$. Note that because of the positivity of the Perron vector, it holds $1>a_1>a_2>\dots>a_{N-1}=v_N>0$.

    As a next step, write
    \begin{align*}
        x_i-\mu&=x_i-x_1-(\mu-x_1)=\sum_{k=1}^{i-1}\delta_k-\sum_{k=1}^{N-1}a_k\delta_k=\sum_{k=1}^{i-1}\delta_k-\sum_{k=1}^{i-1}a_k\delta_k-\sum_{k=i}^{N-1}a_k\delta_k\\&
        =\underbrace{\sum_{k=1}^{i-1}(1-a_k)\delta_k}_{=P_i}-\underbrace{\sum_{k=i}^{N-1}a_k\delta_k}_{=Q_i}
    \end{align*}
    Consequently,
    \begin{align*}
        \|x_i-\mu\|^2&=\|P_i-Q_i\|^2=\|P_i\|^2+\|Q_i\|^2-2\langle P_i,Q_i\rangle\\&
        \le\|P_i\|^2+\|Q_i\|^2+2\langle P_i,Q_i\rangle=\|P_i+Q_i\|^2
    \end{align*}
    for all $i=2,\dots,N-1$, where we have used
    \begin{align*}
        \langle P_i,Q_i\rangle=\sum_{1\le k\le i-1}\sum_{i\le l\le N-1}\underbrace{(1-a_k)}_{>0}\underbrace{a_l}_{>0}\underbrace{\langle\delta_k,\delta_l\rangle}_{\ge0}>0.
    \end{align*}
    In particular, this is the step where we use the acute chain condition~\eqref{eq:acute_chain_condition}.

    Now we see that
    \begin{align}\label{eq:acute_chain_proof_002}
        \|P_i+Q_i\|\le\begin{cases}
            \|\sum_{k=1}^{N-1}a_k\delta_k\|=\|\mu-x_1\| & ,a_{i-1}\ge\frac{1}{2}\\
            \|\sum_{k=1}^{N-1}(1-a_k)\delta_k\|=\|x_N-\mu\| & ,a_{i-1}<\frac{1}{2}.
        \end{cases}
    \end{align}

    Finally, by combining Equations~\eqref{eq:bias_variance_decomp_perron_entry} and~\eqref{eq:acute_chain_proof_002} we arrive at
    \begin{align*}
        \rho v_i-\sigma^2&=\|\mu-x_i\|^2\le\|P_i+Q_i\|^2\le\max\{\|\mu-x_1\|,\|x_N-\mu\|\}\\&
        =\max\{\rho v_1-\sigma^2,\rho v_N-\sigma^2\}
    \end{align*}
    and therefore $v_i\le\max\{v_1,v_N\}$ for all interior coordinates $i=2,\dots,N-1$.
\end{proof}

\subsection{Proof of Lemma~\ref{lem:frechet_derivative_of_kde}}\label{sec:proof_frechet_derivative_of_kde}

\begin{proof}
    Let $h>0$. For every $x\in\mathbb{R}$ we have
    \begin{align*}
        \lim\limits_{\tau\to0}\frac{1}{\tau}(\hat{f}_{h+\tau}(x)-\hat{f}_h(x))&=\frac{1}{n}\sum_{j=1}^{n}\lim\limits_{\tau\to0}\frac{K_{h+\tau}(x-X_j)-K_{h}(x-X_j)}{\tau}\\&
        =\frac{1}{n}\sum_{j=1}^{n}\partial_hK_h(x-X_j)\\&
        =-\frac{1}{nh^2}\sum_{j=1}^{n}\left(K\left(\frac{x-X_j}{h}\right)+\frac{x-X_j}{h}K^{'}\left(\frac{x-X_j}{h}\right)\right)\\&
        =:\partial_h\hat{f}_h(x).
    \end{align*}
    Let $0<|\tau|<h/2$. Using the mean value inequality, we bound
    \begin{align*}
        \left|\frac{K_{h+\tau}(u)-K_h(u)}{\tau}\right|&\le\sup_{s\in(\min\{h,h+\tau\},\max\{h,h+\tau\})}|\partial_sK_s(u)|\\&
        \le\sup_{s\in(h/2,3h/2)}\left|\frac{1}{s^2}L\left(\frac{u}{s}\right)\right|\\&
        \le\frac{4}{h^2}\underbrace{\sup_{a\in(2/3,2)}|L(au/h)|}_{=:G(u/h)}
    \end{align*}
    and
    \begin{align*}
        \left|\frac{\hat{f}_{h+\tau}(x)-\hat{f}_h(x)}{\tau}\right|&\le\frac{1}{n}\sum_{j=1}^{n}\left|\frac{K_{h+\tau}(x-X_j)-K_h(x-X_j)}{\tau}\right|\\&
        \le\frac{4}{nh^2}\sum_{j=1}^{n}G\left(\frac{x-X_j}{h}\right)=:M(x),
    \end{align*}
    where $M\in L^2(\mathbb{R})$ because by the Cauchy-Schwarz inequality $\|M\|_2^2\le 16h^{-3}\|G\|_2^2$ and by the decay assumption on $L$ we get $\|G\|_2^2\le3C^2/(2\alpha-1)<+\infty$. By the dominated convergence theorem
    \begin{align}\label{eq:dom_conv_consequence}
        \left\|\frac{\hat{f}_{h+\tau}-\hat{f}_h}{\tau}-\partial_h\hat{f}_h\right\|_2\to0,\quad\tau\to0.
    \end{align}
    Set $d\Psi(h)[\tau]:=\tau\partial_h\hat{f}_h$. Now it holds
    \begin{align*}
        \hat{f}_{h+\tau}-\hat{f}_h=\tau\partial_h\hat{f}_h+\underbrace{\left(\hat{f}_{h+\tau}-\hat{f}_h-\tau\partial_h\hat{f}_h\right)}_{=:r_h(\tau)},
    \end{align*}
    where the remainder $r_h(\tau)$ fulfills $\|r_h(\tau)\|_2/|\tau|\to0$ as $\tau\to0$ as established in Equation~\eqref{eq:dom_conv_consequence}.

    \textcolor{red}{We must now establish the continuity of $d\Psi:(0,\infty)\to\mathcal{L}(\mathbb{R},L^2(\mathbb{R}))$.}

    Fix a closed nonempty interval $J:=[a,b]\subseteq(0,\infty)$. Applying the bound from above:
    \begin{align*}
        \|r_h(\tau)\|_2&\le|\tau|\sup_{t\in[0,1]}\|d\Psi(h+\tau t)-d\Psi(h)\|_{\text{op}}\\&
        \le|\tau|\underbrace{\sup_{h,s\in J,|h-s|\le|\tau|}\|d\Psi(s)-d\Psi(h)\|_\text{op}}_{=:\mu_{J}(|\tau|)}
    \end{align*}
    for all $h\in[a,b]$ and $\tau\in\mathbb{R}$ with $h+\tau\in[a,b]$. $\mu_{J}(s)=\sup_{x,y\in J,|x-y|\le s}\|d\Psi(x)-d\Psi(y)\|_{\text{op}}$ is known in the literature as the modulus of continuity. Because $d\Psi$ is continuous, $d\Psi:J\to\mathcal{L}(\mathbb{R},L^2(\mathbb{R}))$ is uniformly continuous, and in this case it is well known that $\mu_{J}(s)\to0$ for $s\downarrow0$. Also note that $\|d\Psi(s)-d\Psi(h)\|_{\text{op}}=\|\partial_h\hat{f}_h|_{h=s}-\partial_h\hat{f}_h|_{h=h}\|_2$. Since we have established the continuity of $d\Psi$ before, the proof is finished.
\end{proof}

\subsection{Proof of Proposition~\ref{prop:main_thm}}\label{sec:proof_main_thm}

\begin{proof}
    Set $\epsilon_N:=\sqrt{\lambda_N}\Delta_{r,N}$ and $t_{N,j}:=\epsilon_N(j-i_N)$. Note that $\epsilon_N=t_{N,j+1}-t_{N,j}$ for all $j=1,\dots,N-1$. Fix $A>0$ and set $I_N(A):=\{j=1,\dots,N:|t_{N,j}|\le A\}$. 

    We have $\#I_N(A)\le2A\epsilon_N^{-1}+1$ because $\#I_N(A)\le\#\{j\in\mathbb{Z}:-A\epsilon_N^{-1}+i_N\le j\le A\epsilon_N^{-1}+i_N\}=\lfloor A\epsilon_N^{-1}\rfloor-\lceil-A\epsilon_N^{-1}\rceil+1=2\lfloor A\epsilon_N^{-1}\rfloor+1\le 2A\epsilon_N^{-1}+1$.

    We also set $\delta_{N,j}:=(j-i_N)\Delta_{r,N}=\psi_r(h_j)-\psi_r(h_{i_N})$.

    Split the row sum:
    \begin{align*}
        R_{i_N}^{(N)}=\sum_{j\in I_N(A)}S_{i_Nj}^{(N)}+\underbrace{\sum_{j\notin I_N(A)}S_{i_Nj}^{(N)}}_{=:T_N(A)}.
    \end{align*}

    We have $D_{i_Nj}\ge c(\delta_{N,j})^2$ for some constant $c>0$. To see this, first write
    \begin{align*}
        \left|e^{-\frac{h_j^2\omega^2}{2}}-e^{-\frac{h_{i_N}^2\omega^2}{2}}\right|=\int_{\min\{h_j,h_{i_N}\}}^{\max\{h_j,h_{i_N}\}}t\omega^2e^{-\frac{t^2\omega^2}{2}}\,dt\ge|h_{i_N}-h_j|h_{\min}\omega^2e^{-h_{\max}^2\omega^2/2}
    \end{align*}
    and, furthermore
    \begin{align*}
        D_{i_Nj}&=\frac{1}{4\pi}\int|P_n(\omega)|^2\left|e^{-\frac{h_j^2\omega^2}{2}}-e^{-\frac{h_{i_N}^2\omega^2}{2}}\right|^2\,d\omega\\&
        \ge|h_{i_N}-h_j|^2\frac{h_{\min}^2}{4\pi}\int\underbrace{|P_n(\omega)|^2\omega^4e^{-h_{\max}^2\omega^2}}_{=:g(\omega)}\,d\omega.
    \end{align*}
    Now, $|\delta_{N,j}|=|h_j^r-h_{i_N}^r|/r=\int_{\min\{h_j,h_{i_N}\}}^{\max\{h_j,h_{i_N}\}}s^{r-1}\,ds\le h_{\min}^{r-1}|h_j-h_{i_N}|$. Thus, it follows $D_{i_Nj}\ge c(\delta_{N,j})^2$ with positive constant $c=h_{\min}^{4-2r}(4\pi)^{-1}\int g(\omega)\,d\omega$. It follows $S_{i_Nj}^{(N)}=e^{-\lambda_ND_{i_Nj}}\le e^{-c\lambda_N(\delta_{N,j})^2}=e^{-ct_{N,j}^2}$. Note that the constant is, in fact, positive due to the continuity of the integrand $g(\omega)$ and because $g(\omega)\ne0$ on a disk $\{\omega:|\omega|\le\delta\}\setminus\{0\}$ for some $\delta>0$.

    Continuing
    \begin{align*}
        \epsilon_NT_N(A)&\le\epsilon_N\sum_{1\le j\le N:\epsilon_N|j-i_N|>A}e^{-ct_{N,j}^2}=\epsilon_N\sum_{1-i_N\le k\le N-i_N:|k|>\epsilon_N^{-1}A}e^{-c\epsilon_N^2k^2}\\&
        \le\epsilon_N\sum_{k\in\mathbb{Z}:|k|>\epsilon_N^{-1}A}e^{-c\epsilon_N^2k^2}=2\sum_{k=\lfloor\epsilon_N^{-1}A\rfloor+1}^{\infty}\epsilon_Ne^{-c\epsilon_N^2k^2}\\&
        \le2\sum_{k=\lfloor\epsilon_N^{-1}A\rfloor+1}^{\infty}\int_{(k-1)\epsilon_N}^{k\epsilon_N}e^{-ct^2}\,dt\le2\int_{A-\epsilon_N}^{\infty}e^{-ct^2}\,dt
    \end{align*}
    We see that $\epsilon_NT_N(A)$ is asymptotically bounded and therefore
    \begin{align*}
        \limsup_{N\to\infty}\epsilon_NT_N(A)\le2\int_{A}^{\infty}e^{-ct^2}\,dt\le\frac{2}{A}\int_A^{\infty}te^{-ct^2}\,dt=\frac{e^{-cA^2}}{cA}\to0,\quad A\to\infty.
    \end{align*}
    
    Now, let $j\in I_N(A)$, that is, $j-i_N\le|j-i_N|\le\epsilon_N^{-1}A=\mathbf{O}(N/\sqrt{\lambda_N})=\mathbf{o}(N)$. Thus, by Lemma~\ref{lem:asymptotic_approx_r_linear_spacing} we have $\tau_{N,j}:=h_j-h_{i_N}=h_{i_N}^{1-r}\delta_{N,j}+\mathbf{O}(\delta_{N,j}^2)$ and by Lemma~\ref{lem:frechet_derivative_of_kde} follows
    \begin{align*}
        \hat{f}_{h_j}-\hat{f}_{h_{i_N}}&=\hat{f}_{h_{i_N}+\tau_{N,j}}-\hat{f}_{h_{i_N}}=\tau_{N,j}\partial_{h}\hat{f}_h|_{h=h_{i_N}}+u_{N,j}\\&
        =h_{i_N}^{1-r}\delta_{N,j}\partial_h\hat{f}_h|_{h=h_{i_N}}+r_{N,j},
    \end{align*}
    where $u_{N,j}\in L^2$ is the remainder with $\|u_{N,j}\|_2\le|\tau_{N,j}|\mu_{[h_{\min},h_{\max}]}(|\tau_{N,j}|)$ according to Lemma~\ref{lem:frechet_derivative_of_kde} and $r_{N,j}:=u_{N,j}+\alpha_{N,j}\partial_h\hat{f}_{h_{i_N}}$, where $|\alpha_{N,j}|\le C\delta_{N,j}^2$ for some constant $C>0$ and all $N$ sufficiently large.

    We have $|\tau_{N,j}|\le h_{\max}^{1-r}|\delta_{N,j}|$ and $|\delta_{N,j}|\le A\lambda_N^{-1/2}$ for all $j\in I_N(A)$. Since $h\mapsto\partial_h\hat{f}_h$ is continuous in $L^2(\mathbb{R})$, it follows that $h\mapsto\|\partial_h\hat{f}_h\|_2$ is continuous in $\mathbb{R}$. Thus, $M_0:=\max_{h\in[h_{\min},h_{\max}]}\|\partial_h\hat{f}_h\|_2$ is finite.

    We have $\|r_{N,j}\|_2\le|\delta_{N,j}|a_N(A)$ for some sequence $a_N(A)\to0$ for $N\to\infty$, because
    \begin{align*}
        \|r_{N,j}\|_2&\le\|u_{N,j}\|_2+C\delta_{N,j}^2\|\partial_h\hat{f}_h|_{h=h_{i_N}}\|_2\\&
        \le|\tau_{N,j}|\mu_{[h_{\min},h_{\max}]}(|\tau_{N,j}|)+C\delta_{N,j}^2M_0\\&
        \le|\delta_{N,j}|\left(h_{\max}^{1-r}\mu_{[h_{\min},h_{\max}]}(h_{\max}^{1-r}|\delta_{N,j}|)+CM_0|\delta_{N,j}|\right)\\&
        \le|\delta_{N,j}|\left(\underbrace{h_{\max}^{1-r}\mu_{[h_{\min},h_{\max}]}(h_{\max}^{1-r}A\lambda_N^{-1/2})+CM_0A\lambda_N^{-1/2}}_{=:a_N(A)}\right)
    \end{align*}
    and $a_N(A)\to0$ since $\lambda_N^{-1/2}\to0$ by assumption and $\mu_{[h_{\min},h_{\max}]}(s)\to0$ for $s\downarrow0$ by Lemma~\ref{lem:frechet_derivative_of_kde}.

    Consequently,
    \begin{align*}
        D_{i_Nj}=\frac{1}{2}\|\hat{f}_{h_j}-\hat{f}_{h_{i_N}}\|_2^2=\frac{1}{2}h_{i_N}^{2(1-r)}(\delta_{N,j})^2\|\partial_h\hat{f}_h|_{h=h_{i_N}}\|_2^2+\rho_{N,j}=(\delta_{N,j})^2\hat{Q}_r(h_{i_N})+\rho_{N,j},
    \end{align*}
    where $\rho_{N,j}=h_{i_N}^{1-r}\delta_{N,j}\langle\partial_h\hat{f}_h|_{h=h_{i_N}},r_{N,j}\rangle+(1/2)\|r_{N,j}\|_2^2$ is yet another remainder term for which we have
    \begin{align*}
        |\rho_{N,j}|&\le|\delta_{N,j}|h_{i_N}^{1-r}|\langle\partial_h\hat{f}_{h}|_{h=h_{i_N}},r_{N,j}\rangle|+(1/2)\|r_{N,j}\|_2^2\\&
        \le|\delta_{N,j}|h_{\max}^{1-r}\underbrace{\|\partial_h\hat{f}_h|_{h=h_{i_N}}\|_2}_{\le M_0}\underbrace{\|r_{N,j}\|_2}_{\le|\delta_{N,j}|a_N(A)}+(1/2)\underbrace{\|r_{N,j}\|_2^2}_{\le(\delta_{N,j})^2a_N(A)^2}\\&
        \le(\delta_{N,j})^2\underbrace{\left(h_{\max}^{1-r}M_0a_N(A)+(1/2)a_N(A)^2\right)}_{=:b_N(A)}.
    \end{align*}
    Since $a_N(A)\to0$, we have shown that $\rho_{N,j}=\mathbf{o}(\delta_{N,j}^2)$. Thus,
    \begin{align*}
        S_{i_Nj}=e^{-\lambda_ND_{i_Nj}}=e^{-\lambda_N\delta_{N,j}^2\hat{Q}_r(h_{i_N})}e^{-\lambda_N\rho_{N,j}}.
    \end{align*}
    Define $\eta_{N,j}=e^{-\lambda_N\rho_{N,j}}-1$. It follows $S_{i_Nj}=e^{-\lambda_N(\delta_{N,j})^2\hat{Q}_r(h_{i_N})}(1+\eta_{N,j})$. Now write
    \begin{align*}
        \epsilon_N\sum_{j\in I_N(A)}S_{i_Nj}^{(N)}=\epsilon_N\sum_{j\in I_N(A)}e^{-\lambda_N\delta_{N,j}^2\hat{Q}_r(h_{i_N})}+\epsilon_N\sum_{j\in I_N(A)}e^{-\lambda_N\delta_{N,j}^2\hat{Q}_r(h_{i_N})}\eta_{N,j}.
    \end{align*}
    The local error can be bounded:
    \begin{align*}
        \left|\epsilon_N\sum_{j\in I_N(A)}\underbrace{e^{-\lambda_N\delta_{N,j}^2\hat{Q}_r(h_{i_N})}}_{\le 1}\eta_{N,j}\right|\le\underbrace{\epsilon_N\#I_N(A)}_{\le2A+\epsilon_N}\sup_{j\in I_N(A)}|\eta_{N,j}|\to0,\quad N\to\infty.
    \end{align*}

    We must show that $\sup_{j\in I_N(A)}|\eta_{N,j}|\to0$ for $N\to\infty$. We have $\lambda_N|\rho_{N,j}|\le\lambda_N(\delta_{N,j})^2b_N(A)\le A^2b_N(A)$ for all $j\in I_N(A)$. Consequently, $|\eta_{N,j}|=|e^{-\lambda_N\rho_{N,j}}-1|\le\max\{e^{A^2b_N(A)}-1,1-e^{-A^2b_N(A)}\}\to0$ as $N\to\infty$.

    
    Since $\hat{Q}_r$ is continuous in $h_0$, we have $\hat{Q}_r(h_{i_N})\to\hat{Q}_r(h_0)$ for $N\to\infty$. Therefore
    \begin{align*}
        \left|\epsilon_N\sum_{j\in I_N(A)}e^{-t_{N,j}^2\hat{Q}_r(h_{i_N})}-\epsilon_N\sum_{j\in I_N(A)}e^{-t_{N,j}^2\hat{Q}_r(h_0)}\right|\le(2A+\epsilon_N)A^2|\hat{Q}_r(h_{i_N})-\hat{Q}_r(h_0)|\to0.
    \end{align*}
    We continue working with the remaining non-vanishing term:
    \begin{align*}
        &\left|\epsilon_N\sum_{j\in I_N(A)}e^{-t_{N,j}^2\hat{Q}_r(h_{i_N})}-\int_{-A}^{A}e^{-t^2\hat{Q}_r(h_0)}\,dt\right|\\&\le\underbrace{\left|\epsilon_N\sum_{j\in I_N(A)}e^{-t_{N,j}^2\hat{Q}_r(h_{i_N})}-\epsilon_N\sum_{j\in I_N(A)}e^{-t_{N,j}^2\hat{Q}_r(h_0)}\right|}_{\to0}+\left|\epsilon_N\sum_{j\in I_N(A)}e^{-t_{N,j}^2\hat{Q}_r(h_0)}-\int_{-A}^{A}e^{-t^2\hat{Q}_r(h_0)}\,dt\right|.
    \end{align*}

    In order to show that the second summand above also converges to $0$ for $N\to\infty$ we first rewrite the sum over $I_N(A)$. Before that, set $q=\hat{Q}_r(h_0)$, $g(t):=e^{-qt^2}$, $m_N:=\lfloor A\epsilon_N^{-1}\rfloor$ and $x_k:=\epsilon_N(-m_N+k-1)$. For sufficiently large $N$ we can now write
    \begin{align*}
        \epsilon_N\sum_{j\in I_N(A)}g(t_{N,j})&=\epsilon_N\sum_{j=1,\dots,N:-A\epsilon_N^{-1}+i_N\le j\le A\epsilon_N^{-1}+i_N}g(t_{N,j})\\&
        =\epsilon_N\sum_{k=-\lfloor A\epsilon_N^{-1}\rfloor}^{\lfloor A\epsilon_N^{-1}\rfloor}g(k\epsilon_N)\\&
        =\sum_{k=-m_N}^{m_N}(k\epsilon_N-(k-1)\epsilon_N)g(k\epsilon_N)\\&
        =\sum_{k=1}^{2m_N+1}(x_k-x_{k-1})g(x_k)
    \end{align*}

    We have used $t_{N,1}=\sqrt{\lambda_N}\delta_{N,1}=\sqrt{\lambda_N}r^{-1}(h_{\min}^r-h_{i_N}^r)\to-\infty$ and $t_{N,N}\to+\infty$ for $N\to\infty$ because we assumed that $h_{i_N}\to h_0$ for some interior scale $h_0$ and because $\lambda_N\to\infty$.

    Now $\{x_0,\dots,x_{2m_N+1}\}$ is a partition of the interval $[-\epsilon_Nm_N-\epsilon_N,\epsilon_Nm_N]$.

    Thus,
    \begin{align*}
        \left|\sum_{k=1}^{2m_N+1}(x_k-x_{k-1})g(x_k)-\int_{-\epsilon_Nm_N-\epsilon_N}^{\epsilon_Nm_N}g(x)\,dx\right|&\le\sum_{k=1}^{2m_N+1}\int_{x_{k-1}}^{x_k}|g(x_k)-g(x)|\,dx\\&
        \le2A\hat{Q}_r(h_0)(2A\epsilon_N+\epsilon_N^2)\to0,\quad N\to\infty,
    \end{align*}
    because $\epsilon_N\to0$ since we have assumed $\lambda_N/N^2\to0$ and
    \begin{align*}
        \left|\int_{-\epsilon_Nm_N-\epsilon_N}^{\epsilon_Nm_N}g(t)\,dt-\int_{-A}^{A}g(t)\,dt\right|\to0,\quad N\to\infty,
    \end{align*}
    because $-\epsilon_N(m_N+1)\to-A$ and $\epsilon_Nm_N\to A$.

    We have now shown that the aforementioned second summand also converges to zero.

    Finally,
    \begin{align*}
        L_A&=\lim_{N\to\infty}\epsilon_N\sum_{j\in I_N(A)}S_{i_Nj}^{(N)}\le\lim_{N\to\infty}\epsilon_N\sum_{j\in I_N(A)}S_{i_Nj}^{(N)}+\liminf_{N\to\infty}\epsilon_NT_N(A)\\&
        \le\liminf_{N\to\infty}\epsilon_NR_{i_N}^{(N)}\le\limsup_{N\to\infty}\epsilon_NR_{i_N}^{(N)}\le\lim_{N\to\infty}\epsilon_N\sum_{j\in I_N(A)}S_{i_Nj}^{(N)}+\limsup_{N\to\infty}\epsilon_NT_N(A)\\&
        \le L_A+\frac{e^{-cA^2}}{cA}
    \end{align*}
    
    Letting $A\to\infty$, we see that the limit of $\epsilon_NR_{i_N}^{(N)}$ as $N\to\infty$ exists and is given by
    \begin{align*}
        \lim_{N\to\infty}\epsilon_NR_{i_N}^{(N)}=\lim_{A\to\infty}L_A=\int_{-\infty}^{\infty}e^{-t^2\hat{Q}_r(h_0)}\,dt=\sqrt{\frac{\pi}{\hat{Q}_r(h_0)}}.
    \end{align*}
    We must make sure that $\hat{Q}_r(h_0)>0$. This follows from previously established results inside this proof. For all sufficiently large $N$ we have $0<c\le\delta_{N,i_N+1}^{-2}D_{i_N,i_N+1}=\hat{Q}_r(h_{i_N})+\delta_{N,i_N+1}^{-2}\rho_{N,i_N+1}$ and taking the limit gives $\hat{Q}_r(h_0)\ge c>0$.
\end{proof}